\documentclass{ics}

% import packages here
% \usepackage{foo,bar}

\begin{document}
\title{Sample Texts for Symposium on Innovations in Computer Science}
\author{Ruini Xue$^{1}$ \and Wenguang Chen$^{2}$}
\address{1. Tsinghua University Beijing 100084, 2. Insititute of High Performance Computing Beijing 100084}
\email{xueruini@gmail.com, cwg@tsinghua.edu.cn}

\begin{abstract} 
  This template follows \LaTeXe{} conventions, and your \TeX{} and \LaTeX{}
  experiences should work here. The only exception is: call
  \texttt{$\backslash$maketitle} after abstract and keywords! And please
  separate key words with semicolon. Finally, please first update \LaTeX{}
  packages if you encounter strange problems.
\end{abstract}

\keywords{word1; word2; word3; and word4}

%% Important: do call \maketitle after abstract and keywords!!!
\maketitle


\section{Introduction}
\label{sec:introduction}
The Institute for Computer Science (ITCS) is a new institute headed by Professor
Andrew Chi-Chih Yao at Tsinghua University. Its mission is to nurture and
develop the talents of students and researchers and to create knowledge for the
advancement of computer science. It also aims to become one of the leading
research centers on theoretical computer science in the world. The Institute
already has an elite graduate student body and enjoys frequent visits by
world-renowned scientists all year round. With a vibrant research environment
and impressive research outputs, ITCS provides an ideal setting for fruitful
interactions and collaborations among its members.

\subsection{References Example}
\label{sec:ref-example}
Xue and Chen explored the feasibility of user-level file system in HPC
environment~\cite{xue08}. Donald~\onlinecite{tex} wrote the masterpiece. Some
other references~\cite{xue09,Chafik94}

\subsection{Details}
\label{sec:details}
The Institute for Computer Science (ITCS) is a new institute headed by Professor
Andrew Chi-Chih Yao at Tsinghua University. Its mission is to nurture and
develop the talents of students and researchers and to create knowledge for the
advancement of computer science. It also aims to become one of the leading
research centers on theoretical computer science in the world. The Institute
already has an elite graduate student body and enjoys frequent visits by
world-renowned scientists all year round. With a vibrant research environment
and impressive research outputs, ITCS provides an ideal setting for fruitful
interactions and collaborations among its members.

\subsubsection{More Details}
\label{sec:more-details}
The Institute for Computer Science (ITCS) is a new institute headed by Professor
Andrew Chi-Chih Yao at Tsinghua University. Its mission is to nurture and
develop the talents of students and researchers and to create knowledge for the
advancement of computer science. It also aims to become one of the leading
research centers on theoretical computer science in the world. The Institute
already has an elite graduate student body and enjoys frequent visits by
world-renowned scientists all year round. With a vibrant research environment
and impressive research outputs, ITCS provides an ideal setting for fruitful
interactions and collaborations among its members.

\subsection{Equation Example}
\label{sec:equation-example}
Please check Eq.~\eqref{eq:1}. DO read \textsf{amsmath} document before you
input complicated equations.

\begin{equation}
  \label{eq:1}
  E=mc^{2}
\end{equation}

\paragraph{More More Details}
\label{sec:more-more-details}
The Institute for Computer Science (ITCS) is a new institute headed by Professor
Andrew Chi-Chih Yao at Tsinghua University. Its mission is to nurture and
develop the talents of students and researchers and to create knowledge for the
advancement of computer science. It also aims to become one of the leading
research centers on theoretical computer science in the world. The Institute
already has an elite graduate student body and enjoys frequent visits by
world-renowned scientists all year round. With a vibrant research environment
and impressive research outputs, ITCS provides an ideal setting for fruitful
interactions and collaborations among its members.

\subparagraph{The Details}
\label{sec:details-1}
The Institute for Computer Science (ITCS) is a new institute headed by Professor
Andrew Chi-Chih Yao at Tsinghua University. Its mission is to nurture and
develop the talents of students and researchers and to create knowledge for the
advancement of computer science. It also aims to become one of the leading
research centers on theoretical computer science in the world. The Institute
already has an elite graduate student body and enjoys frequent visits by
world-renowned scientists all year round. With a vibrant research environment
and impressive research outputs, ITCS provides an ideal setting for fruitful
interactions and collaborations among its members.

\subsection{Table Example}
\label{sec:more-more-details-1}
Please check Table~\ref{tab:example}.

\begin{table}[h]
  \centering
  \caption{Table Example}
  \label{tab:example}
  \begin{tabular}{|l|l|l|}
    \hline
    first & second & third \\
    \hline
    first & second & third \\
    \hline
    first & second & third \\
    \hline
  \end{tabular}
\end{table}

\section{Design}
\label{sec:design}
The Institute for Computer Science (ITCS) is a new institute headed by Professor
Andrew Chi-Chih Yao at Tsinghua
University~\footnote{\url{http://www.tsinghua.edu.cn}}. Its mission is to
nurture and develop the talents of students and researchers and to create
knowledge for the advancement of computer science. It also aims to become one of
the leading research centers on theoretical computer science in the world. The
Institute already has an elite graduate student body and enjoys frequent visits
by world-renowned scientists all year round. With a vibrant research environment
and impressive research outputs, ITCS provides an ideal setting for fruitful
interactions and collaborations among its members.

\subsection{Architecture}
\label{sec:architecture}
The Institute for Computer Science (ITCS) is a new institute headed by Professor
Andrew Chi-Chih Yao at Tsinghua University. Its mission is to nurture and
develop the talents of students and researchers and to create knowledge for the
advancement of computer science. It also aims to become one of the leading
research centers on theoretical computer science in the world. The Institute
already has an elite graduate student body and enjoys frequent visits by
world-renowned scientists all year round. With a vibrant research environment
and impressive research outputs, ITCS provides an ideal setting for fruitful
interactions and collaborations among its members.

\subsubsection{Implementation}
\label{sec:implementation}
The Institute for Computer Science (ITCS) is a new institute headed by Professor
Andrew Chi-Chih Yao at Tsinghua University. Its mission is to nurture and
develop the talents of students and researchers and to create knowledge for the
advancement of computer science. It also aims to become one of the leading
research centers on theoretical computer science in the world. The Institute
already has an elite graduate student body and enjoys frequent visits by
world-renowned scientists all year round. With a vibrant research environment
and impressive research outputs, ITCS provides an ideal setting for fruitful
interactions and collaborations among its members.

\subsection{Figure Example}
\label{sec:ref-figure}
Please check Figure~\ref{fig:example}. The figure may not appear the place where
your code locates, this is because figures and tables are float objects in
\LaTeX, which means that they are can be moved around to generate better
typesetting results.

\begin{figure}
  \centering
  \includegraphics[width=.8\linewidth]{thu-whole-logo}
  \caption{Figure Example}
  \label{fig:example}
\end{figure}

\section*{Acknowledegment}
\label{sec:acknowledegment}
This research is supported by NSF Grant 123456.

% References
% 
% \bibliographystyle{unsrt}
% \bibliography{bibfile}

\begin{thebibliography}{99}
%\bibitem{journal} Author1, Author2. Paper Title[J]. Journal Name, Year, Volume(Number):pages.
\bibitem{xue08} Ruini X, Wenguang C, Weimin Z. CprFS: A User-level File System
  to Support Consistent File States for Checkpoint and Restart. 22nd ACM
  International Conference on Supercomputing (ICS'08), 114-223, June 7-12, 2008.

%\bibitem{book} Author1, Author2. Book Title[M]. editions. Publish address:publisher,year.
\bibitem{tex} Knuth D~E. The {\TeX} Book. 15th ed. Reading, MA: Addison-Wesley
  Publishing Company, 1989

\bibitem{xue09} Ruini X, Xuezheng L, Ming W, Zhenyu G, Wenguang C, Weimin Z,
  Zheng Z, Geoffery M~V. MPIWiz: Subgroup Reproducible Replay of MPI
  Applications. 14th ACM SIGPLAN Symposium on Principles and Practice of
  Parallel Programming (PPoPP'09), 251-260, February 14-18, 2009.

\bibitem{Chafik94} {Chafik El Idrissi} M, Roney A, Frigon C, et~al. Measurements
  of total kinetic-energy released to the {$N=2$} dissociation limit of {H}$_2$
  --- evidence of the dissociation of very high vibrational {R}ydberg states of
  {H}$_2$ by doubly-excited states. Chemical Physics Letters, 1994,
  224(10):260--266

%\bibitem{onlinedata} Author. Report Title[EB/OL]. [publish date]. the link
%\bibitem{website} Web Site Name. the link

\end{thebibliography}
\end{document}